(1) 在Euler–Bernoulli小挠度与线弹性范围内，横向挠度的稳定方程是 \(EI\,y''''+P\,y''=0\)。两端铰支要求 \(y(0)=y(L)=0\)，同时弯矩 \(M_b=-EI y''\) 在端点为零，故 \(y''(0)=y''(L)=0\)；不能把铰支误作夹固端的零转角条件。
(2) 令 \(y(x)=A\sin(kx)\)，代入得 \(EI k^4-Pk^2=0\)，非零模满足 \(P=EI k^2\)。边界使 \(k_n=n\pi/L\)，因此 \(P_n=n^2\pi^2EI/L^2\)，\(n=1,2,\ldots\)；最小的第一屈曲临界载荷为 \(P_{\rm cr}=\pi^2EI/L^2\)。
(3) 直杆附近的二次势能为 \(\delta^2\Pi=\frac12\int_0^L[EI(y'')^2-P(y')^2]dx\)。对满足 \(y(0)=y(L)=0\) 的扰动有Wirtinger型不等式 \(\int_0^L(y'')^2dx\ge(\pi/L)^2\int_0^L(y')^2dx\)（还用到 \(y'\) 的零平均）。所以若 \(P<P_{\rm cr}\)，\(\delta^2\Pi>0\)，直杆线性稳定。
(4) 由 \(P_{\rm cr}=\pi^2EI/L^2\) 可见长度加倍使临界载荷下降到四分之一，而截面弯曲惯性矩 \(I\) 加倍使临界载荷加倍。该推导要求细长杆、材料未屈服、轴向载荷无偏心且屈曲幅度足够小；临界点以上的真实后屈曲形状须采用非线性几何。
